\chapter{Att kontrollera ett påstående}
\label{app:verifiera}

Det här materialet gör, som allt undervisningsmaterial, en rad
påståenden.  De flesta går att pröva vid datorn: att ett stavfel i ett
nyckelnamn upptäcks först när raden körs ser du själv.  Men några av
påståendena handlar om vad forskningen och fältet säger, och dem kan
ingen körning avgöra:
\begin{enumerate}
  \item att man bör föredra komposition framför arv --- att rådet är
    etablerat i objektorienterad design och kommer från Gamma, Helm,
    Johnson och Vlissides (\cref{sec:designprinciper,sec:komposition-arv});
  \item att inkapsling är en etablerad princip
    (\cref{sec:designprinciper});
  \item att en klass med ett enda ansvar följer principen \emph{SRP},
    som är Robert C. Martins (\cref{sec:designprinciper});
  \item att det är bra att inte upprepa sig i ett program --- principen
    \emph{DRY} --- och att den kommer från Hunt och Thomas
    (\cref{sec:arv-medborgare});
  \item att den enklaste lösning som fungerar är att föredra ---
    \emph{KISS} --- och att principen tillskrivs Kelly Johnson
    (\cref{sec:varor}).
\end{enumerate}
Hur vet man om sådant stämmer?  Samma fråga möter dig varje gång du läser
en lärobok, en webbsida eller ett svar från en språkmodell --- och varje
gång du själv ska skriva en rapport.  Den här bilagan sammanfattar
arbetsgången, och \cref{tab:paastaenden} visar frågorna och
svaren.  Påståendena är redan prövade i tidigare föreläsningar, och
tabellen anger var; det första prövas i föreläsningen \emph{Klasser
och objekt}, bilaga~D, där rådet först ges.  Arbetsgången beskrivs
utförligare i bilagan med samma namn i föreläsningen \emph{Algoritmiskt
tänkande}.

\begin{table}[htbp]
  \centering
  \caption{Påståendena som frågor, och var svaren finns.}
  \label{tab:paastaenden}
  \begin{tabular}{@{}lp{0.42\linewidth}p{0.3\linewidth}@{}}
    \toprule
    Bilaga & Fråga & Svar \\
    \midrule
    --- & Bör man föredra komposition framför arv, och kommer rådet
      från Gamma med flera? & Redan belagt i \emph{Klasser och objekt},
      bilaga~D \\
    --- & Är inkapsling en etablerad princip? & Redan belagt i
      \emph{Klasser och objekt}, bilaga~C \\
    --- & Är SRP Robert C. Martins princip, och är den etablerad? &
      Redan belagt i \emph{Funktioner}, bilaga~B \\
    --- & Är det bra att inte upprepa sig, och kommer DRY från Hunt och
      Thomas? & Redan belagt i \emph{Algoritmiskt tänkande}, bilaga~E \\
    --- & Tillskrivs KISS Kelly Johnson, och är enkelhet ett etablerat
      ideal? & Redan belagt i \emph{Funktioner}, bilaga~C \\
    \bottomrule
  \end{tabular}
\end{table}

\section{Gör påståendet till en fråga}

Ett påstående går att kontrollera först när det är tydligt vad som skulle
kunna visa att det är \emph{fel}.  Formulera det därför som en fråga med
ett svar.  \enquote{Bör man föredra komposition framför arv?} kan
besvaras med ja eller nej; \enquote{komposition är bra} kan det inte,
eftersom \enquote{bra} är ett omdöme och inte ett faktum.  Ofta döljer
sig flera påståenden i ett: \enquote{ett etablerat råd som kommer från
Gamma med flera} är två frågor --- \emph{ursprung} och \emph{etablering}
--- och de kräver olika slags källor; och \emph{om rådet är riktigt} är
en tredje, empirisk fråga.

\section{Välj rätt sorts källa}

Olika frågor besvaras av olika källor.
\begin{description}
  \item[Uppslagsverk] för vad ett ord \emph{betyder}.  För vad ett
    programspråk \emph{gör} är språkets egen dokumentation motsvarande
    källa.
  \item[Primärkällan] för var en idé \emph{kommer ifrån}, och för ett
    specifikt resultat.  Vill man veta vad Gamma med flera faktiskt
    skrev läser man deras bok, inte en webbsida som återger den ---
    i det här fallet fanns boken i författarens eget bibliotek.
  \item[Översikter] (\foreignlanguage{english}{reviews}) för om något är
    \emph{etablerat}, eller för vad forskningen \emph{sammantaget}
    säger.
  \item[Enskilda studier] för ett specifikt resultat --- med förbehållet
    att en enskild studie bara visar att något observerats en gång, i en
    viss population.
\end{description}
Var man hittar dem: universitetsbibliotekets söktjänst, och databaser som
Scopus, Web of Science, IEEE Xplore, DBLP (datalogi) och OpenAlex
(öppen).  I kursens bilagor används kommandoradsverktyget
\texttt{scholar}\footnote{Paketet \texttt{scholarcli} på PyPI:
  \url{https://pypi.org/project/scholarcli/}.}, som ställer samma fråga
till flera databaser på en gång och sparar allt som hände i en
\emph{session}.

\section{Sök så att du kan ha fel}

Den som söker på namn hon redan känner till hittar det hon väntar sig.
Sök därför \emph{ämnesdrivet} --- på begrepp, inte författare.  Några
praktiska regler:
\begin{itemize}
  \item Håll frågorna korta.  Flera korta frågor slår en lång.
  \item Ställ \emph{samma} frågor i alla databaser.
  \item Sök också efter \emph{invändningar}: lägg till ord som
    \enquote{critique}, \enquote{limitations} eller \enquote{benefits}
    av det som påståendet avråder från.  Ett påstående som överlevt en
    motsökning är värt mer än ett som aldrig prövats, och en invändning
    som hittas blir ett \emph{förbehåll} i svaret, inte ett nederlag.
    Bilaga~D i föreläsningen \emph{Klasser och objekt}, om komposition
    och arv, är just ett sådant fall.
  \item Räkna med att databasernas rankning missar somligt.  Då hämtas
    källan som \emph{känt dokument}: man vet att den finns och slår upp
    den direkt.  Det är i sin ordning, bara det sägs.
\end{itemize}

\section{Läs i fulltext, och skriv ner hur}

En träff i en databas är inte ett belägg.  Sammanfattningen
(\foreignlanguage{english}{abstract}) räcker inte heller: den säger vad
författarna vill ha sagt, inte vad de visat.  Läs hela texten och leta
upp \emph{det ställe} som styrker påståendet --- ett ordagrant citat med
sidnummer.  Räkna aldrig en källa du inte läst.  Går fulltexten inte
att nå skriver man det, och nöjer sig med det sammanfattningen faktiskt
säger; så är det med flera av källorna i föreläsningen \emph{Klasser
och objekt}, bilaga~D.

Skriv sedan ner hur du gick till väga, så att någon annan kan följa
spåret.  I det här materialet står den anteckningen som en kommentar
ovanför varje post i litteraturlistans källfil, med fasta fält:
\begin{description}
  \item[\texttt{CLAIM}] vilket påstående källan ska styrka;
  \item[\texttt{FOUND-VIA}] hur den hittades --- sökfråga och databas,
    eller \enquote{känt dokument};
  \item[\texttt{PICKED}] varför just den, framför andra träffar;
  \item[\texttt{QUOTE}] det ordagranna citatet, med sida;
  \item[\texttt{VERIFIED}] att fulltexten lästs, och var;
  \item[\texttt{COUNTER}] motsökningen och vad den gav;
  \item[\texttt{DATE}] när.
\end{description}
Det tar ett par minuter per källa, och det är det enda som gör det
möjligt att om ett år svara på frågan \enquote{hur vet du det?}

\section{Dra slutsatsen --- med förbehåll}

Svaret skrivs ut tillsammans med det som talar emot.  Bilaga~D i
föreläsningen \emph{Klasser och objekt} svarar \enquote{ja, med
förbehåll} om komposition och arv: rådet är
etablerat och arvets kostnader är belagda, men de kontrollerade
experimenten om arv och underhållbarhet är oeniga, och arv används
flitigt och förnuftigt i verklig kod.  Förbehållen är inte svagheter i
svaret.  De \emph{är} svaret, och de avgör hur påståendet används i
materialet: tumregeln i \cref{sec:komposition-arv} säger
\enquote{föredra}, inte \enquote{undvik}.  Varje bilaga som redovisar
en sökning slutar därför med ett avsnitt \emph{Fortsatt arbete}: vad
just den sökningen lämnade ogjort, och vad som skulle behöva undersökas
för ett säkrare svar.

\section{En anmärkning om språkmodeller}

Sökningen om komposition och arv gav hundratals träffar.  En
språkmodell användes för att \emph{sortera} dem --- hör träffen till ämnet över huvud taget,
och åt vilket håll pekar den? --- så att gallringen kunde granskas.
Modellen sorterade, men varje källa som citeras ska läsas och väljas av
en människa, och även sorteringen ska bekräftas av en människa.
Modellens klassning står därför med sin konfidens i den bilagans
tabell, så att
den kan ifrågasättas.  Använd gärna språkmodeller för att hitta och
sortera; låt dem aldrig vara det sista ledet.
